% !TEX root = ../projekt.tex
% =============================================================================
% STAV: návrhy a nápady (bod 4 zadania). Implementované je iba to, čo je
% v tabuľke označené. Ostatné je plán, nie výsledok.
% =============================================================================
\chapter{Návrh variantov jednotky RSS}
\label{chap:varianty}

Analýza ukázala, že hash nie je úzkym miestom. Varianty sa preto zameriavajú na to,
\emph{ako sa z~hashu alebo zo stavu kanálov určí kanál}. Podľa zadania sa delia do dvoch skupín.

\begin{table}[ht]
  \centering
  \small
  \caption{Prehľad variantov a ich aktuálny stav.}
  \label{tab:varianty}
  \begin{tabular}{l l l}
    \toprule
    Varianta & Afinita toku & Stav \\
    \midrule
    A1 Iná hashovacia funkcia     & áno & vyhodnotené, bez prínosu \\
    A2 Symetria                   & áno & vyhodnotené, má cenu \\
    A3 Väčšia RETA                & áno & návrh \\
    A4 Úprava hashu (digest)      & áno & návrh \\
    A5 RETA podľa záťaže          & áno & návrh \\
    A6 Orákulum JSQ               & áno & hotové, dolná hranica \\
    \midrule
    B1 Skupiny susedných kanálov  & nie & návrh \\
    B2 Flowlety                   & nie & návrh \\
    B3 Dva hashe a výber v~podmnožine & nie & návrh \\
    B4 Rozdelenie iba veľkých tokov & nie & návrh \\
    \bottomrule
  \end{tabular}
\end{table}

\section{Varianty so zachovaním afinity toku}

\subsection*{A1: Iná hashovacia funkcia}
Vyhodnotené v~sekciách~\ref{sec:hash-nezalezi} až~\ref{sec:statistika}. Medzi kvalitnými funkciami
sa prínos nepreukázal. Slabšia funkcia môže pri vstupe s~nízkou entropiou výsledok zhoršiť
(sekcia~\ref{sec:dalsie-vzorky}).

\subsection*{A2: Symetria}
Vyhodnotené v~sekcii~\ref{sec:cena-symetrie}. Je to funkčná požiadavka s~merateľnou cenou,
nie optimalizácia.

\subsection*{A3: Väčšia tabuľka RETA}
Odstraňuje kvantizáciu pri počtoch kanálov, ktoré nedelia veľkosť tabuľky. Pri 512 položkách
klesne teoretické preťaženie dokonalého hashu z~37,5 na 9,4\,\textperthousand{} ($n=20$)
a z~50,0 na 12,5\,\textperthousand{} ($n=40$). \TODO{Zistiť cenu v~pamäti FPGA.}

\subsection*{A4: Úprava hashu pred výberom položky}
Medzi hash a RETA sa vloží lacná operácia, napríklad zmiešanie horných bitov do spodných
($h \oplus (h \gg 16)$). \emph{Očakávanie:} vzhľadom na výsledky A1 skôr malý prínos, je to však
najlacnejšia zmena.

\subsection*{A5: Plnenie RETA podľa záťaže}
Položky RETA sa medzi oknami presúvajú z~preťažených kanálov na menej zaťažené. Hardvér meria
záťaž položiek a tabuľku preprogramuje. Je to realizovateľná aproximácia orákula.
\TODO{Otvorené: ako často meniť tabuľku a čo s~tokmi, ktoré sa pri zmene presunú
(krátkodobé preusporiadanie paketov).}

\emph{Dopad na simulátor:} kanál sa dnes určí raz pri štarte toku. Pre A5 sa musí pre tok uložiť
položka RETA a kanál určovať až pri skladaní okna.

\subsection*{A6: Orákulum JSQ}
Popísané v~sekcii~\ref{sec:jsq}. Nie je to návrh pre hardvér, ale referencia, s~ktorou sa
porovnávajú ostatné varianty.

\section{Varianty rozdeľujúce tok na podmnožinu kanálov}

Tieto varianty útočia priamo na príčinu z~sekcie~\ref{sec:slony}: veľký tok už nemusí celý
zaťažiť jeden kanál. \textbf{Cenou je možné preusporiadanie paketov toku}, preto k~nim treba
pridať novú metriku. \TODO{Navrhnúť metriku preusporiadania a overiť s~vedúcim, koľko kanálov
v~skupine je prijateľných pre aplikácie.}

\subsection*{B1: Skupiny susedných kanálov}
Hash vyberie skupinu susedných kanálov a pakety toku sa v~nej striedajú (\emph{round robin}).
Pôvodný nápad: skupiny po 2 kanáloch pri 8 kanáloch a po 3 až 6 pri 128 kanáloch, lebo hardvér
vie susedné kanály prepojiť.

\subsection*{B2: Flowlety}
Tok sa rozdelí v~medzerách medzi paketmi väčších než zvolený prah a každá časť sa hashuje zvlášť.
\emph{Obmedzenie simulátora:} pakety sú v~toku rozložené rovnomerne, takže skutočné medzery
nevidí. Na túto variantu by bola potrebná paketová stopa. \TODO{Zvážiť, či ju vôbec zaradiť.}

\subsection*{B3: Dva hashe a výber v~podmnožine}
Hash určí podmnožinu $k$ kanálov a v~nej sa pre paket vyberie kanál druhým mechanizmom:
striedaním, druhým hashom alebo podľa najmenšej záťaže.

\subsection*{B4: Rozdelenie iba veľkých tokov}
Malé toky zostávajú v~jednom kanáli, rozdelia sa iba toky, ktoré počas behu prekročia prah.
Väčšina tokov sa teda nepreusporiada a záťaž veľkých tokov sa rozloží.

\emph{Dopad na simulátor pre skupinu B:} kanál sa musí vyberať pre každý paket, nie raz pre tok.

\TODO{Vybrať 3--4 varianty, ktoré sa naozaj implementujú. Návrh: A3, A5, B1, B4.}
